% Part: first-order-logic
% Chapter: axiomatic-deduction
% Section: axioms-rules-quantifiers

\documentclass[../../../include/open-logic-section]{subfiles}

\begin{document}

\olfileid{fol}{axd}{qua}

\olsection{Axioma's en regels voor kwantoren}

\begin{defn}[Axioma's voor kwantoren]
De \emph{axioma's} voor kwantoren, de tekens voor `voor alle' en `er
bestaat', zijn alle invullingen van de volgende twee vormen:
\begin{align}
\ollabel{ax:q1} & \lforall[x][!B] \lif !B(t), \\
\ollabel{ax:q2} & !B(t) \lif \lexists[x][!B].
\end{align}
Hierbij mag $t$ elke gesloten term zijn, dus elke term zonder
variabelen.
\end{defn}

\begin{defn}[Regels voor kwantoren]
\item Kijk eerst naar $!B \lif !A(a)$. Staat deze formule al in de
  !!{derivation}, en komt $a$ niet voor in~$\Gamma$ en ook niet in~$!B$? Dan
  telt $!B \lif \lforall[x][!A(x)]$ als een correcte stap.
\item Kijk nu naar $!A(a) \lif !B$. Staat deze formule al in de
  !!{derivation}, en komt $a$ niet voor in~$\Gamma$ en ook niet in~$!B$? Dan
  telt $\lexists[x][!A(x)] \lif !B$ als een correcte stap.
\end{defn}

Voor elk van deze twee regels schrijven we kortweg ``\QR.''

\end{document}
